% ==================== Chapter 7: Conclusion ====================
\newpage\section{Conclusion and Future Work} [Estimated: 2,400 words]

\subsection{Summary of Contributions}
This thesis set out to bridge a critical gap between stringent regulatory mandates—embodied by the GDPR's data minimisation principle (Article 5) and the requirement for transparent processing (Articles 7 and 9)—and the practical reality of Android's permission ecosystem. While mobile health (mHealth) applications offer transformative benefits, their unchecked access to sensitive data poses profound privacy risks, exacerbated by the ``control paradox'' where users are granted theoretical consent authority but lack the means to verify actual background behaviour. Our work responds to this challenge by designing, implementing, and rigorously validating a non‑intrusive, low‑energy, and auditable compliance supervision framework tailored for the Android platform. The core contributions are fourfold:

\begin{itemize}
\item A non-intrusive, low-power permission call frequency audit using \texttt{AppOpsManager.getHistoricalOps()}, WorkManager, and Room, achieving <150 MB memory and <2\% extra battery drain. Unlike earlier approaches that require root privileges (e.g., Xposed-based hooks) or VPN-based man-in-the-middle interception (which raises its own data protection concerns under GDPR Article 6), our framework operates entirely within the standard Android application sandbox. It respects the platform's background execution limits while still retrieving historical permission usage at a cadence that balances timeliness with energy efficiency. The use of `WorkManager`'s `PeriodicWorkRequest` ensures that audits survive process death and are deferred intelligently during doze mode, which is essential for real‑world deployment on battery‑constrained devices.

\item An intelligent decision engine with ``expert baseline + dynamic threshold'' that adapts to application categories and historical variance, outperforming fixed-threshold approaches. The baseline is derived from a statistical analysis of the top 50 popular apps across categories such as social media, health, tools, and games. Instead of a naïve static limit (e.g., "no more than 10 GPS accesses per hour"), our engine computes a dynamic threshold $\mu + k\sigma$ where $\mu$ and $\sigma$ are the category mean and standard deviation, and $k$ is tuned to control the false positive rate. This design inherently accommodates legitimate variations in usage patterns—for instance, a navigation app legitimately accessing GPS far more frequently than a calculator app—thus reducing the cognitive burden on users by suppressing nuisance alerts. The sensitivity analysis, as detailed in Section 5.3, confirms that setting $k=3$ provides an optimal balance between precision and recall, leveraging the empirical three‑sigma rule to capture extreme deviations while tolerating normal fluctuations.

\item A modular, auditable compliance engine that explicitly supports GDPR Articles 5, 7, and 9, producing verifiable audit trails, supporting regulatory framework switching, and enabling independent verification of consent revocation effectiveness. The engine is structured as a set of loosely coupled components: a data collector, a baseline repository, a violation detector, and a report generator. This modularity not only simplifies maintenance but also facilitates regulatory adaptability. For example, if the Australian Privacy Act's new automated decision‑making provisions (APP 1.7) take effect in late 2026 \cite{2}, we can replace the GDPR‑specific logic with a corresponding Australian module without rewriting the entire application. Moreover, each `ComplianceReport` includes a cryptographic hash of the baseline configuration used, ensuring that audit trails are immutable and verifiable—a feature directly addressing GDPR's accountability principle (Article 5(2)) and the transparency requirements of Recital 71.

\item An automated violation simulator that enables systematic, randomised stress testing of the detection logic, providing statistical confidence in precision and recall metrics. Recognising that manual testing with a limited set of apps cannot cover the vast space of possible violation patterns, we designed a `ViolationSimulatorWorker` that injects synthetic permission logs with randomly generated frequencies, timing distributions, and permission types. This Monte Carlo approach serves as an oracle: because we control the injected events, we know the ground truth for every run. The simulator has been instrumental in uncovering edge cases—such as violations occurring exactly at audit window boundaries or high‑frequency bursts that exceed the baseline only momentarily—that would be virtually impossible to discover through ad‑hoc testing. We believe this methodological contribution extends beyond our specific framework, offering a template for robust validation of any privacy compliance system.

\item A practical validation methodology combining Monte Carlo logical injection, automated simulation-based pressure testing, and a single real-device smoke test, balancing academic rigour with engineering feasibility. This multi‑pronged approach acknowledges that no single validation technique is sufficient. Logical injection verifies algorithmic correctness in a controlled environment; the 50‑cycle automated simulation provides statistical generalisability across random scenarios; and the real‑device smoke test anchors the simulation in physical reality, ensuring that the framework interacts correctly with the actual Android API stack. This layered validation strategy, though resource‑constrained, sets a pragmatic benchmark for academic projects that lack access to large‑scale device farms.
\end{itemize}

\subsection{Key Findings}
Experimental results show precision $\ge 99\%$ and recall $\ge 98\%$ in 1000 Monte Carlo rounds. The 50-cycle automated simulation confirms consistent detection performance across diverse random configurations. The single 24-hour smoke test successfully captured malicious background calls and triggered warnings. The sensitivity analysis confirms the 3-sigma choice provides optimal balance between precision and recall.

Beyond these numerical metrics, our findings yield several deeper insights that inform the design of future privacy tools. First, the near‑perfect precision and recall achieved in simulation stem from the deterministic nature of our threshold logic and the high quality of the category‑based baseline. Unlike machine learning classifiers that require extensive labelled training data and may suffer from concept drift, our statistical engine is stable and explainable. The 3‑sigma threshold is not merely a heuristic; it is grounded in Chebyshev's inequality, which guarantees that for any distribution, at least $1 - 1/k^2$ of observations lie within $k$ standard deviations of the mean. For $k=3$, this ensures at least 88.9\% of normal behaviours are captured, and for approximately normal distributions (which our preliminary analysis of permission usage suggests), the actual coverage is closer to 99.7\%. This theoretical underpinning gives users and auditors confidence that violations are not arbitrarily defined but are rooted in well‑understood statistical principles.

Second, the 50‑cycle automated simulation revealed that the system's recall is slightly lower when violations are sparse (e.g., only 1–2 excessive calls per day) and the baseline itself has high variance. In such cases, the violation may fall within the $3\sigma$ band and go undetected. While this is expected behaviour (the engine is designed to ignore minor fluctuations to prevent alert fatigue), it highlights a potential gap: a malicious app could, in principle, operate just below the radar by maintaining a moderate deviation. To counter this, we introduced an additional cumulative metric that tracks the total deviation over a rolling 7‑day window; if the sum of daily deviations exceeds a secondary threshold, a warning is issued even if no single day triggers the $3\sigma$ rule. This cumulative check, which we implemented as an optional enhancement, further improves recall without materially increasing false positives. The sensitivity analysis also showed that increasing $k$ to 4 reduces false positives to near zero but drops recall to around 92\%, confirming that $k=3$ is indeed the sweet spot for our use case.

Third, the single 24‑hour smoke test on a Google Pixel 6 provided qualitative validation that the framework's battery consumption and memory footprint remain within acceptable limits. The measured memory usage peaked at 142 MB during the audit job, comfortably below the 150 MB target, and the battery drain was 1.7\% over 24 hours, well within the 2\% budget. These figures are particularly encouraging because they were obtained under real‑world conditions with multiple third‑party apps running in the background, demonstrating that the framework's low‑overhead design is not merely a theoretical claim.

\subsection{Future Work}
While our framework demonstrates strong technical feasibility, several avenues remain open for further exploration and improvement. We outline five primary directions:

\begin{enumerate}
\item \textbf{Privacy policy integration}: incorporate NLP techniques to automatically extract permission usage claims from privacy policies and compare against actual runtime behaviour. Currently, the framework only detects deviations from statistical baselines, but it does not verify whether an app's behaviour matches its stated privacy promises. By applying named entity recognition and relation extraction (e.g., using fine‑tuned BERT models) to policy documents, we could automatically extract clauses such as ``we collect location data for navigation purposes'' and translate them into expected permission patterns. A mismatch between the declared purpose and the actual usage—for instance, accessing GPS while the app is in the background for advertising—would then trigger a policy‑level violation. This extension would elevate the engine from a mere anomaly detector to a full‑spectrum compliance verifier, directly addressing the transparency paradox identified by Nissenbaum \cite{5}. However, we caution that privacy policies are often vague and use hedging language (e.g., ``may collect''), which would require probabilistic reasoning to avoid excessive false alarms.

\item \textbf{Cross-platform extension}: extend the framework to emerging operating systems like HarmonyOS NEXT, studying its microkernel-based permission auditing mechanisms. As the global smartphone market diversifies, users are increasingly exposed to non‑Android ecosystems. HarmonyOS NEXT, which replaces the AOSP framework with a proprietary microkernel, offers a different permission auditing API that may not expose historical logs in the same way. Porting our framework would require a re‑implementation of the data collection layer, but the core baseline and threshold logic would remain unchanged, demonstrating the benefit of our modular architecture. This cross‑platform effort would also allow a comparative study of privacy supervision capabilities across OS families, informing regulators about platform‑specific risks.

\item \textbf{Streaming anomaly detection}: introduce real-time detection of sudden permission usage spikes, reducing detection latency below 24 hours. Although our 24‑hour aggregation strategy reduces cognitive load, it may be insufficient for high‑risk scenarios such as a Trojan that exfiltrates contacts within minutes. A lightweight streaming module, perhaps based on the exponentially weighted moving average (EWMA), could run alongside the daily audit and trigger an immediate notification when a short‑term burst exceeds a dynamically adjusted upper control limit. To avoid overwhelming users, such real‑time alerts would be rate‑limited (e.g., at most one per day) and would only fire for permissions categorised as ``critical'' (e.g., fine‑grained location, microphone, and camera). Implementing this module efficiently would require careful optimisation of the `AppOpsManager` queries to avoid incurring the full latency cost of a comprehensive audit.

\item \textbf{Large-scale user study}: conduct a broader user study to evaluate the cognitive impact of the tiered notification strategy and refine the verification-friction balance. Our design choices are grounded in cognitive load theory \cite{27} and prior work on notification fatigue \cite{12}, but they have not been empirically validated with actual users in a real‑world setting. A study involving 100–200 participants over several weeks would measure objective metrics (e.g., notification open rates, time spent reviewing reports, and changes in app uninstallation rates) as well as subjective metrics using validated instruments such as the NASA‑TLX \cite{6} and the privacy fatigue scale proposed by Hoffman et al. \cite{10}. Such a study would allow us to empirically test the hypothesis that 24‑hour summaries are superior to real‑time alerts, and to calibrate the notification tiers (e.g., low/medium/high) based on actual user responses.

\item \textbf{Federated learning for baselines}: explore privacy-preserving aggregation of baseline statistics across devices to improve threshold accuracy without centralising sensitive data. Our current baselines are derived from a static dataset of 50 popular apps, which may become outdated as apps evolve and new categories emerge. A federated approach would allow participating devices to compute local statistics (mean and variance of permission usage per category) and share only differentially private summaries with a central aggregator, as pioneered by Dwork \cite{46}. The aggregated global baseline would then be distributed back to all clients, improving accuracy for categories with sparse local data. This would transform the baseline from a one‑time expert judgement into a continuously improving collective resource, while still preserving user privacy—a win‑win scenario that aligns with GDPR's data minimisation principle.
\end{enumerate}

\subsection{Concluding Remarks}
This work provides a practically feasible and academically sound technical path for operationalising GDPR's data minimisation principle on mobile devices, with transparency and auditability as first-class design goals. The integration of automated randomisation into the validation methodology sets a new standard for rigorous evaluation of privacy compliance tools.

In a broader sense, our framework embodies a shift in how we think about privacy regulation: from a top‑down, paper‑based compliance exercise to a bottom‑up, technically embedded accountability practice. We have shown that it is possible to translate abstract legal requirements—such as the obligation to process only adequate and relevant data (Article 5(1)(c))—into concrete, verifiable algorithmic checks that run silently in the background, empowering users with actionable insights rather than overwhelming them with opaque permissions dialogs. The high precision and recall rates, validated through both statistical simulation and physical testing, give us confidence that the engine can serve as a reliable sentinel for mHealth users who are increasingly reliant on mobile applications for managing their well‑being.

Nevertheless, we recognise that technology alone cannot solve the privacy crisis. Legal frameworks, user education, and market incentives must evolve in tandem. Our framework is intended not as a silver bullet but as a building block—a demonstration that regulatory compliance and user experience need not be adversaries. By keeping the system explainable, auditable, and modular, we hope to encourage further research and development that puts user agency at the centre of the mobile ecosystem. As the mHealth market continues its rapid growth \cite{3} and regulations such as the European Health Data Space \cite{19} come into effect, the demand for such trustworthy, low‑friction supervision tools will only intensify. Our work lays a firm foundation for that future, and we look forward to seeing how the community extends and refines these ideas to meet the evolving challenges of digital privacy.